\documentclass[11pt]{article}
\usepackage[a4paper,margin=1in]{geometry}
\usepackage{amsmath,amssymb,mathtools,bm}
\usepackage{enumitem,booktabs,array}
\usepackage{tcolorbox}
\usepackage{hyperref}
\usepackage[protrusion=true,expansion=false]{microtype}
\hypersetup{colorlinks=true,linkcolor=blue,urlcolor=blue,citecolor=blue}
\setlist[itemize]{topsep=2pt,itemsep=2pt}
\setlist[enumerate]{topsep=2pt,itemsep=2pt}
\newtcolorbox{keybox}{colback=blue!3!white,colframe=blue!40!black,boxrule=0.4pt,arc=1mm}
\newtcolorbox{alertbox}{colback=red!3!white,colframe=red!40!black,boxrule=0.4pt,arc=1mm}
\newtcolorbox{answerbox}{colback=green!3!white,colframe=green!40!black,boxrule=0.4pt,arc=1mm}
\newtcolorbox{drillbox}{colback=yellow!5!white,colframe=orange!60!black,boxrule=0.4pt,arc=1mm}
\newcommand{\EV}{\mathbb{E}}
\newcommand{\Var}{\mathrm{Var}}
\newcommand{\Cov}{\mathrm{Cov}}
\newcommand{\blank}[1]{\underline{\hspace{#1}}}

\title{W05-02: Egan, Hortaçsu \& Matvos (2017, AER) \\
\large ``Deposit Competition and Financial Fragility: Evidence from the US Banking Sector'' --- Active Recall Workbook}
\author{Banking \& Risk Management --- Exam Prep}
\date{\today}
\begin{document}\maketitle

\begin{keybox}
\textbf{Paper in one sentence.} EHM estimate a structural model of US deposit competition combining a depositor discrete-choice demand model with a global-games bank-solvency model: insured and uninsured depositors differ in sensitivity to bank risk, banks face multiple equilibria (``good''/``bad'') in deposit rates, and the model delivers quantitative counterfactuals --- e.g., removing deposit insurance triples fragility-linked welfare costs.

\textbf{Placement.} The structural-IO take on bank runs and deposit markets. Pairs with Iyer--Puri (W05-01) on depositor behaviour and with DSS (W02-03) on the market power of deposits.
\end{keybox}

\section*{Round 1: Complete Walk-Through}

\subsection*{1.1 Framework}
Depositors choose banks based on rate, branch network, and \emph{perceived failure probability}. Banks set deposit rates given depositor demand and a solvency condition that can be characterised by a global-games threshold (Goldstein--Pauzner style): failure occurs if enough uninsured depositors run.

\subsection*{1.2 Demand side}
Two depositor types:
\begin{itemize}
\item \emph{Insured} (retail, below \$250k): rate-sensitive, low sensitivity to bank failure risk.
\item \emph{Uninsured} (large corporates, above \$250k): strongly sensitive to failure probability.
\end{itemize}
Discrete-choice logit with bank-region-time fixed effects. Elasticity of uninsured-deposit share to perceived failure probability is large and negative.

\subsection*{1.3 Supply side and equilibrium}
Each bank chooses a deposit rate trading off higher rates (more deposits) against higher funding cost and higher failure probability (endogenous through leverage). With global games, the model can support multiple equilibria: a ``good'' low-rate, low-failure equilibrium and a ``bad'' high-rate, high-failure equilibrium.

\subsection*{1.4 Estimation}
Structural estimation on US bank Call-Report and RateWatch data 2002--2013. Identify (i) demand parameters from cross-bank variation in rates/shares, (ii) supply parameters from rate-setting responses to shocks.

\subsection*{1.5 Headline results}
\begin{itemize}
\item Uninsured depositors' demand is highly elastic to perceived bank risk; insured much less.
\item Multiple equilibria appear empirically for realistic parameter regions.
\item Counterfactual: removing deposit insurance triples the welfare costs associated with bank failures (run-driven losses).
\item Counterfactual: a regulator-imposed rate cap narrows the ``bad'' equilibrium region and stabilises the system.
\end{itemize}

\subsection*{1.6 Caveats}
\begin{alertbox}
\begin{itemize}
\item Structural model is stylised; many simplifications (homogeneous banks within region, etc.).
\item Counterfactuals are only as valid as the model; large extrapolations require caution.
\item Global-games selection is a theoretical device, not directly observed in data; EHM use the model's implications.
\end{itemize}
\end{alertbox}

\section*{Round 2: Level 1 Fill-in}
\begin{enumerate}
\item Model blocks: \blank{2cm}-choice depositor demand $+$ \blank{2cm}-games bank solvency.
\item Insured vs.\ uninsured: differ in sensitivity to \blank{2cm} risk.
\item Equilibrium multiplicity: \blank{1cm} rate-low failure and \blank{1cm} rate-high failure.
\item Data span: \blank{1cm}--\blank{1cm}.
\item Counterfactual: remove insurance $\Rightarrow$ welfare cost \blank{1cm}.
\end{enumerate}

\section*{Round 3: Level 2 Prompts}
\begin{enumerate}
\item Write the depositor logit demand and explain identification of the failure-probability coefficient.
\item Why can multiple equilibria arise in the deposit market?
\item Give the intuition for the ``triple welfare cost'' counterfactual.
\item How does EHM's framework connect to DSS's deposits channel?
\end{enumerate}

\section*{Round 4: Minimal Scaffolding}
From a blank page: (i) demand block, (ii) supply block, (iii) multiple-equilibria result, (iv) headline counterfactual, (v) two caveats.

\section*{Round 5: Blank Page}
Essay: EHM's structural view of financial fragility --- what does the IO lens add to Diamond--Dybvig?

\vspace{6pt}
\begin{drillbox}
\textbf{Speed Drill.} (1) Two modelling blocks. (2) Insured vs.\ uninsured sensitivity. (3) Multiple-equilibria statement. (4) Counterfactual without insurance. (5) Data period.
\end{drillbox}

\section*{Quick Reference \& Answer Key}
\begin{answerbox}
\begin{itemize}
\item \textbf{Model.} Logit demand $+$ global-games solvency.
\item \textbf{Depositor types.} Insured low-sensitivity, uninsured high-sensitivity to risk.
\item \textbf{Equilibria.} Multiple: good / bad high-rate.
\item \textbf{Counterfactual.} No insurance $\Rightarrow$ $\sim 3\times$ fragility welfare cost.
\end{itemize}
\end{answerbox}

\begin{keybox}
\textbf{30-second exam pitch.} Egan, Hortaçsu \& Matvos (2017) bring IO structural methods to the bank-run literature. Depositor demand is logit over rates and failure-probability perceptions, with insured depositors essentially insensitive to failure risk and uninsured depositors highly responsive. Banks set deposit rates under a global-games solvency condition, which in realistic parameter regions gives rise to multiple equilibria: a good low-rate/low-failure equilibrium and a bad high-rate/high-failure equilibrium. Structural estimation on 2002--2013 US Call-Report and RateWatch data recovers the key elasticities. Counterfactuals quantify policy: eliminating deposit insurance roughly triples the welfare cost of fragility; rate caps stabilise by shrinking the bad-equilibrium region. The paper is the canonical demonstration that deposit competition and financial fragility must be analysed jointly.
\end{keybox}

\end{document}
